\section{Distinction Graph Machines}
\label{sec:algorithms}

The tree model is a useful finite special case, but it is not the most natural runtime data structure. A tree makes every decision through one irreversible path. A more flexible object is a graph of concepts: local memories live at prototypes, and distinctions live on edges between prototypes that have been confused. This section defines that model.

\subsection{State and prediction}

Let $\phi_\theta:\calX\to\R^d$ be an encoder. A Distinction Graph Machine has runtime state
\begin{align*}
\calS_t=\left(\{c_j,M_j,n_j\}_{j=1}^{N_t},G_t\right).
\end{align*}
Here $c_j\in\R^d$ is a prototype, $M_j$ is a local predictor or memory, $n_j$ is a usage count, and $G_t$ is a graph on the prototypes. An edge $(i,j)$ stores a distinction separating the two concepts, typically
\begin{align*}
u_{ij}=\frac{c_i-c_j}{\norm{c_i-c_j}_2},\qquad b_{ij}=\frac{1}{2}u_{ij}^\top(c_i+c_j).
\end{align*}
The edge distinction is
\begin{align*}
g_{ij}(h)=\one\{u_{ij}^\top h>b_{ij}\}.
\end{align*}

Given $h=\phi_\theta(x)$, the machine computes distances $d_j(h)$ to prototypes and a sparse neighborhood
\begin{align*}
N_k(h)\subseteq\{1,\ldots,N_t\}.
\end{align*}
It then forms routing weights
\begin{align*}
\alpha_j(h)=\frac{\exp(-d_j(h)/T)}{\sum_{r\in N_k(h)}\exp(-d_r(h)/T)}
\qquad j\in N_k(h),
\end{align*}
and predicts by local mixture
\begin{align*}
\hat y=\sum_{j\in N_k(h)}\alpha_j(h)M_j(h).
\end{align*}
Hard nearest-prototype routing is the limit $T\to0$.

\begin{algorithm}[H]
\caption{\textsc{Predict} for a Distinction Graph Machine}
\label{alg:dgm-predict}
\begin{algorithmic}[1]
\State \textbf{Input:} state $\calS_t=(\{c_j,M_j,n_j\}_{j=1}^{N_t},G_t)$, encoder $\phi_\theta$, example $x$
\State $h\gets\phi_\theta(x)$
\State $N_k(h)\gets$ the $k$ nearest prototypes to $h$
\State compute weights $\alpha_j(h)$ for $j\in N_k(h)$
\State \Return $\sum_{j\in N_k(h)}\alpha_j(h)M_j(h)$
\end{algorithmic}
\end{algorithm}

\subsection{Refine or repair on a graph}

After observing consequence $y$, DGM first asks whether an existing concept can absorb the evidence. Let $j^\star$ be the nearest or highest-weight prototype. Define a local inconsistency score
\begin{align*}
\rho_t=L(y,M_{j^\star}(h)).
\end{align*}
If $\rho_t$ is small, the information structure is treated as adequate and the machine repairs the local state:
\begin{align*}
M_{j^\star}\leftarrow\operatorname{Repair}(M_{j^\star},h,y),
\qquad
c_{j^\star}\leftarrow c_{j^\star}+\eta_c(h-c_{j^\star}),
\qquad
n_{j^\star}\leftarrow n_{j^\star}+1.
\end{align*}
If $\rho_t$ is large and the neighborhood contains incompatible consequences, the machine refines by creating a new prototype
\begin{align*}
c_{N_t+1}=h,\qquad M_{N_t+1}=\operatorname{InitMemory}(h,y),
\end{align*}
and adds distinction edges from the new prototype to the confused neighbors.

\begin{algorithm}[H]
\caption{\textsc{Observe} for a Distinction Graph Machine}
\label{alg:dgm-observe}
\begin{algorithmic}[1]
\State \textbf{Input:} state $\calS_t$, labeled consequence $(x,y)$
\State $h\gets\phi_\theta(x)$
\If{$N_t=0$}
\State create prototype $c_1=h$ with memory $\operatorname{InitMemory}(h,y)$
\State \Return updated state
\EndIf
\State $j^\star\gets\argmin_j d_j(h)$
\State $\rho\gets L(y,M_{j^\star}(h))$
\If{$\rho\le\tau_{\mathrm{repair}}$}
\State repair $M_{j^\star}$ and move $c_{j^\star}$ toward $h$
\Else
\State create a new prototype $c_{N_t+1}=h$
\State initialize $M_{N_t+1}$ from $(h,y)$
\For{confused neighbor $i$ of $h$}
\State add edge $(i,N_t+1)$ with distinction $(u_{i,N_t+1},b_{i,N_t+1})$
\EndFor
\EndIf
\State \Return updated state
\end{algorithmic}
\end{algorithm}

\begin{theorem}[Prototype distinction separates its endpoints, informal version of Theorem~\ref{thm:prototype-edge:formal}]
\label{thm:prototype-edge:informal}
If $c_i\ne c_j$ and $u_{ij},b_{ij}$ are defined as above, then
\begin{align*}
u_{ij}^\top c_i-b_{ij}=\frac{1}{2}\norm{c_i-c_j}_2,
\qquad
u_{ij}^\top c_j-b_{ij}=-\frac{1}{2}\norm{c_i-c_j}_2.
\end{align*}
Thus each edge distinction separates the two prototypes that caused it.
\end{theorem}

\subsection{Why graph memory is different from a tree}

A tree stores distinctions as a hierarchy. A mistake near the root changes every downstream routing decision. A graph stores distinctions locally between concepts. Creating a new concept does not rewrite old paths; it adds a local alternative and edges explaining why the new concept differs from its neighbors.

A binary distinction tree is recovered by restricting $G_t$ to a rooted binary hierarchy and forcing every input to follow a single path. DGM is therefore the parent data structure; trees are a sparse, hard-routing special case.

\subsection{Local memories}

The local memory $M_j$ can be chosen according to the task.

\begin{itemize}
\item {\bf Classification.} $M_j$ stores label counts or exemplars associated with prototype $j$.
\item {\bf Regression.} $M_j$ stores a local mean, local linear model, or ridge regressor.
\item {\bf Associative repair.} $M_j$ stores a Projection Memory matrix $W_j$ and updates it by
\begin{align*}
W_{j,t+1}=W_{j,t}+\frac{\eta}{1+\eta\norm{k_t}_2^2}(v_t-W_{j,t}k_t)k_t^\top.
\end{align*}
\end{itemize}

The conceptual rule is unchanged: if the prototype is adequate, repair its local state; if evidence is incompatible with the current prototypes, create a new distinction.

\subsection{Language modeling instance}
\label{sec:language-dgm}

For language modeling, the world state is a prefix and the consequence is the next observed token. Let a frozen or slowly trained backbone produce hidden states
\begin{align*}
h_t=\phi_\theta(x_{\le t}).
\end{align*}
DGM is inserted as a runtime memory over hidden states. Each prototype $c_j$ represents a recurring context type: a citation pattern, a local syntax regime, a document-specific entity usage, or a repeated task format. The local memory $M_j$ predicts a correction to the next-token distribution.

A scalable instance stores a target embedding mean
\begin{align*}
r_j=\frac{1}{n_j}\sum_{s\in C_j}E[x_{s+1}],
\end{align*}
where $E[y]\in\R^{d_e}$ is the output embedding of token $y$. For a new prefix, DGM retrieves prototypes $N_k(h_t)$ and forms the correction
\begin{align*}
\ell_{\mathrm{DGM}}(y\mid h_t)=\sum_{j\in N_k(h_t)}\alpha_j(h_t)E[y]^\top r_j.
\end{align*}
The final logits are
\begin{align*}
\ell(y\mid x_{\le t})=\ell_\theta(y\mid x_{\le t})+\lambda_t\ell_{\mathrm{DGM}}(y\mid h_t).
\end{align*}
Here $\lambda_t$ is a learned or fixed trust weight. After the next token $x_{t+1}$ is observed, the machine updates by \textsc{Observe} with consequence $y=x_{t+1}$. If an existing prototype is repaired, its embedding memory can be updated by
\begin{align*}
r_j^+=\frac{n_jr_j+E[x_{t+1}]}{n_j+1},
\qquad
n_j^+=n_j+1.
\end{align*}
If the observed next token is surprising under the nearest prototype and the context is not well explained by existing prototypes, DGM creates a new prototype at $h_t$ and initializes
\begin{align*}
c_{N_t+1}=h_t,
\qquad
r_{N_t+1}=E[x_{t+1}],
\qquad
n_{N_t+1}=1.
\end{align*}

This gives a direct test-time learning protocol for language modeling. During streaming prediction, document continuation, user correction, tool feedback, or supervised next-token evaluation, the next observed token is genuine evidence and can refine or repair the runtime memory. During unconstrained generation, the model should not update from its own sampled token by default, because self-generated samples are actions, not external consequences.

\begin{algorithm}[H]
\caption{\textsc{ObserveNextToken} for DGM language modeling}
\label{alg:dgm-lm}
\begin{algorithmic}[1]
\State \textbf{Input:} prefix $x_{\le t}$, observed token $x_{t+1}$, DGM state $\calS_t$
\State $h_t\gets\phi_\theta(x_{\le t})$
\State retrieve $N_k(h_t)$ and compute logits $\ell_\theta+\lambda_t\ell_{\mathrm{DGM}}$
\State $j^\star\gets$ nearest or highest-weight prototype
\State $\rho\gets-\log p(x_{t+1}\mid x_{\le t})$
\If{$\rho\le\tau_{\mathrm{repair}}$}
\State update the local token memory at $j^\star$
\Else
\State create a new prototype at $h_t$ with local memory initialized by $E[x_{t+1}]$
\State add distinction edges to confused neighboring prototypes
\EndIf
\State \Return updated state
\end{algorithmic}
\end{algorithm}

\subsection{Complexity}

With exact nearest-neighbor search, prediction costs
\begin{align*}
O(N_td+kd+C_{\mathrm{local}}),
\end{align*}
where $C_{\mathrm{local}}$ is the cost of evaluating the selected local memories. With approximate nearest-neighbor indexing, the $N_td$ term can be replaced by the query cost of the index. Creating a prototype costs $O(d)$ plus the cost of initializing its memory. Adding $q$ distinction edges costs $O(qd)$. If each prototype stores an embedding mean in $\R^{d_e}$, runtime memory is
\begin{align*}
O(N_t(d+d_e)+|E_t|d).
\end{align*}

\begin{theorem}[DGM data-structure bound, informal version of Theorem~\ref{thm:dpm-complexity:formal}]
\label{thm:dpm-complexity:informal}
For exact search over $N_t$ prototypes in $\R^d$, one DGM prediction costs $O(N_td+kd+C_{\mathrm{local}})$. One observation that repairs a selected prototype costs this plus the local repair cost. One observation that creates a new prototype and $q$ distinction edges costs an additional $O(d+qd)$. With embedding-mean memories, the runtime state uses $O(N_t(d+d_e)+|E_t|d)$ memory.
\end{theorem}
